% 03-kien-truc.tex — Chương 3: Kiến trúc hệ thống đề xuất
\chapter{Kiến trúc hệ thống đề xuất}
\label{ch:kientruc}

Trên nền tảng lý thuyết đã trình bày ở chương \ref{ch:coso}, chương này phân tích bài toán
cụ thể của VNPost và đề xuất kiến trúc trợ lý ảo. Mục \ref{sec:baigtoan} làm rõ đặc thù
bài toán; mục \ref{sec:yeucau} xác định yêu cầu; mục \ref{sec:multiagent} trình bày kiến
trúc đa tác tử; mục \ref{sec:pipeline-rag} và mục \ref{sec:pipeline-ft} mô tả pipeline
\rag{} và pipeline \ft; mục \ref{sec:kien-truc-lai} đề xuất kiến trúc lai; và mục
\ref{sec:hatang} trình bày hạ tầng và bảo mật.

\section{Đặc thù bài toán tại VNPost}
\label{sec:baigtoan}

\subsection{Đặc điểm dữ liệu nội bộ}

Dữ liệu nội bộ của VNPost có một số đặc điểm ảnh hưởng trực tiếp đến thiết kế hệ thống:

\begin{itemize}
  \item \textbf{Đa dạng định dạng:} tài liệu tồn tại dưới dạng văn bản số (docx, pdf, txt)
  và cả bản scan. Điều này đòi hỏi bước nhận dạng ký tự quang học và chuẩn hóa về một định
  dạng thống nhất trước khi lập chỉ mục.
  \item \textbf{Cấu trúc phức tạp:} nhiều văn bản chứa bảng biểu (biểu phí, thời gian toàn
  trình, bảng phân loại hàng hóa), tiêu đề nhiều cấp và số hiệu văn bản. Việc chia đoạn
  phải bảo toàn bảng và ngữ cảnh tiêu đề.
  \item \textbf{Ngôn ngữ chuyên ngành:} nhiều thuật ngữ, từ viết tắt và số hiệu đặc thù
  (ví dụ số quyết định, mã dịch vụ), đòi hỏi cấu hình truy hồi thưa phù hợp với tiếng Việt.
  \item \textbf{Tính động:} biểu phí và quy định thay đổi theo thời gian, nên tri thức cần
  được cập nhật mà không phải huấn luyện lại mô hình.
  \item \textbf{Tính bảo mật:} tài liệu nội bộ và dữ liệu khách hàng không được phép rời
  khỏi hạ tầng của Tổng công ty.
\end{itemize}

\subsection{Các nhóm nghiệp vụ}

Hệ thống phục vụ ba nhóm nghiệp vụ chính: (i) hỏi đáp trên tài liệu văn thư và sản phẩm
-- dịch vụ; (ii) tra cứu tình trạng bưu gửi; và (iii) tra cứu thông tin bưu cục. Ngoài ra
còn các nghiệp vụ mở rộng như tra cứu chất lượng phát và tra cứu tài chính -- kế toán.

\section{Yêu cầu chức năng và phi chức năng}
\label{sec:yeucau}

\subsection{Yêu cầu chức năng}

\begin{itemize}
  \item Trả lời câu hỏi tự nhiên dựa trên tài liệu nội bộ, có trích dẫn nguồn.
  \item Ngữ cảnh hóa câu trả lời theo lịch sử hội thoại nhiều lượt.
  \item Tra cứu tình trạng bưu gửi và thông tin bưu cục qua kết nối hệ thống nghiệp vụ.
  \item Định tuyến yêu cầu đến tác tử chuyên biệt phù hợp và hợp nhất kết quả.
  \item Hỗ trợ tạo phiếu yêu cầu (ticket) khi cần chuyển sang hỗ trợ thủ công.
  \item Gợi ý câu hỏi mở đầu phù hợp với nhu cầu người dùng.
\end{itemize}

\subsection{Yêu cầu phi chức năng}

\begin{itemize}
  \item \textbf{Chính xác và có căn cứ:} câu trả lời phải bám sát tài liệu, hạn chế tối đa
  ảo giác.
  \item \textbf{Bảo mật:} toàn bộ mô hình và dữ liệu vận hành cục bộ, không gửi ra ngoài.
  \item \textbf{Độ trễ chấp nhận được:} phản hồi trong vài giây cho câu hỏi thông thường.
  \item \textbf{Khả năng mở rộng:} dễ thêm tác tử, kho tài liệu và mô hình mới.
  \item \textbf{Quan sát được:} có vết, log và chỉ số để giám sát và cải tiến.
  \item \textbf{Kiểm soát truy cập:} phân quyền theo vai trò đối với dữ liệu nhạy cảm.
\end{itemize}

\section{Kiến trúc tổng thể đa tác tử}
\label{sec:multiagent}

Đề tài đề xuất kiến trúc \emph{đa tác tử} gồm một tác tử điều phối trung tâm (host
orchestrator) và các tác tử chuyên biệt. Tác tử điều phối chịu trách nhiệm hiểu yêu cầu,
duy trì ngữ cảnh hội thoại, quyết định định tuyến và hợp nhất câu trả lời. Các tác tử
chuyên biệt đảm nhận từng miền nghiệp vụ hẹp và được gọi qua giao thức A2A; mỗi tác tử
truy cập công cụ của mình qua giao thức MCP.

\begin{figure}[htbp]
\centering
\begin{tikzpicture}[
  node distance=9mm,
  host/.style={draw=vnpostdark,fill=vnpostgreen,text=white,rounded corners,minimum width=3.4cm,minimum height=1.1cm,align=center,font=\small\bfseries},
  agent/.style={draw=vnpostgreen,fill=boxblue,rounded corners,minimum width=3.1cm,minimum height=1cm,align=center,font=\small},
  sys/.style={draw=vnpostdark,fill=white,rounded corners,minimum width=2.4cm,minimum height=0.9cm,align=center,font=\scriptsize},
  arrow/.style={-{Latex[length=2mm]},thick,vnpostdark},
  biarrow/.style={{Latex[length=2mm]}-{Latex[length=2mm]},thick,vnpostdark}
]
\node[host] (host) {Tác tử điều phối (Host)};
\node[agent,below left=16mm and 30mm of host] (spdv) {Tác tử RAG\\sản phẩm -- dịch vụ};
\node[agent,below=16mm of host] (buugui) {Tác tử tra cứu\\bưu gửi};
\node[agent,below right=16mm and 30mm of host] (buucuc) {Tác tử tra cứu\\bưu cục};
\node[sys,below=of spdv] (qdrant) {Qdrant + reranker};
\node[sys,below=of buugui] (track) {Hệ thống tracking};
\node[sys,below=of buucuc] (pg) {CSDL bưu cục};
\draw[biarrow] (host) -- (spdv);
\draw[biarrow] (host) -- (buugui);
\draw[biarrow] (host) -- (buucuc);
\draw[arrow] (spdv) -- (qdrant);
\draw[arrow] (buugui) -- (track);
\draw[arrow] (buucuc) -- (pg);
\node[font=\scriptsize,above=2mm of host] {Người dùng $\leftrightarrow$ Host (giao thức hội thoại)};
\node[font=\scriptsize,below=1mm of host] {A2A};
\end{tikzpicture}
\caption{Kiến trúc đa tác tử: host điều phối các tác tử chuyên biệt qua A2A.}
\label{fig:multiagent}
\end{figure}

Ưu điểm của kiến trúc này là tách biệt trách nhiệm rõ ràng: mỗi tác tử có prompt, công cụ
và kho dữ liệu riêng, dễ kiểm thử và nâng cấp độc lập. Tác tử điều phối chỉ cần biết
``khi nào gọi ai'', nhờ đó giảm tải nhận thức và tăng độ ổn định.

\section{Pipeline RAG chi tiết}
\label{sec:pipeline-rag}

\subsection{Pha lập chỉ mục}

Pha lập chỉ mục biến tài liệu thô thành các đoạn vector trong Qdrant qua các bước:

\begin{enumerate}
  \item \textbf{Thu thập và chuẩn hóa:} tài liệu được thu thập, với các bản scan được xử lý
  OCR bằng mô hình nhận dạng (ví dụ \texttt{olmocr-2-7b} hoặc \texttt{lightonocr2-1b}).
  \item \textbf{Chuyển sang Markdown:} tài liệu được chuyển thành văn bản có cấu trúc để bảo
  toàn tiêu đề và bảng.
  \item \textbf{Trích xuất metadata:} mỗi tài liệu được gắn metadata như tiêu đề, số hiệu
  văn bản, ngày ban hành và tóm tắt do mô hình ngôn ngữ sinh ra.
  \item \textbf{Chia đoạn:} văn bản được chia thành các đoạn nhận biết cấu trúc, với độ dài
  tối đa và độ chồng lấn hợp lý, bảo toàn bảng biểu.
  \item \textbf{Nhúng và lưu trữ:} mỗi đoạn được nhúng thành vector bằng mô hình nhúng và
  lưu vào Qdrant, cùng vector thưa cho BM25.
\end{enumerate}

\subsection{Pha truy hồi và sinh}

Khi có câu hỏi, hệ thống thực hiện:

\begin{enumerate}
  \item \textbf{Nhúng câu hỏi} bằng cùng mô hình nhúng đã dùng khi lập chỉ mục.
  \item \textbf{Truy hồi lai:} chạy song song truy hồi thưa BM25 và truy hồi dày, mỗi nhánh
  lấy top-$k_1$ ứng viên (ví dụ 20).
  \item \textbf{Hợp ứng viên:} gộp kết quả hai nhánh và loại trùng.
  \item \textbf{Xếp hạng lại:} dùng cross-encoder xếp hạng lại toàn bộ ứng viên và giữ
  top-$k_2$ (ví dụ 5) làm ngữ cảnh.
  \item \textbf{Sinh câu trả lời:} mô hình ngôn ngữ nhận câu hỏi, ngữ cảnh truy hồi và lịch
  sử hội thoại để sinh câu trả lời có trích nguồn.
\end{enumerate}

\begin{figure}[htbp]
\centering
\begin{tikzpicture}[
  node distance=7mm,
  proc/.style={draw=vnpostgreen,fill=boxblue,rounded corners,minimum width=3.6cm,minimum height=9mm,align=center,font=\small},
  arrow/.style={-{Latex[length=2mm]},thick,vnpostdark}
]
\node[proc] (doc) {Tài liệu nội bộ (docx/pdf/txt/scan)};
\node[proc,below=of doc] (ocr) {OCR + chuẩn hóa Markdown};
\node[proc,below=of ocr] (meta) {Trích xuất metadata (title, số hiệu, ngày, tóm tắt)};
\node[proc,below=of meta] (chunk) {Chia đoạn nhận biết cấu trúc (bảo toàn bảng)};
\node[proc,below=of chunk] (emb) {Nhúng vector Qwen3 + vector thưa BM25};
\node[proc,below=of emb] (store) {Lưu Qdrant (dense COSINE + sparse IDF)};
\node[proc,below=of store] (ret) {Truy hồi lai top-20/20 $\rightarrow$ hợp $\rightarrow$ rerank $\rightarrow$ top-5};
\node[proc,below=of ret] (gen) {LLM sinh câu trả lời có trích nguồn};
\draw[arrow] (doc) -- (ocr);
\draw[arrow] (ocr) -- (meta);
\draw[arrow] (meta) -- (chunk);
\draw[arrow] (chunk) -- (emb);
\draw[arrow] (emb) -- (store);
\draw[arrow] (store) -- (ret);
\draw[arrow] (ret) -- (gen);
\end{tikzpicture}
\caption{Pipeline RAG đầy đủ từ lập chỉ mục đến sinh câu trả lời.}
\label{fig:pipeline-rag}
\end{figure}

\subsection{Tham số hóa và cấu hình}

Các tham số chính gồm số ứng viên mỗi nhánh truy hồi ($k_1$), số đoạn sau xếp hạng lại
($k_2$) và ngưỡng điểm tối thiểu để loại bỏ ngữ cảnh không liên quan. Cấu hình thực tế được
trình bày ở chương \ref{ch:trienkhai}. Cơ chế ngưỡng điểm giúp hệ thống từ chối trả lời
khi không tìm thấy bằng chứng phù hợp, thay vì ``bịa'' câu trả lời.

\section{Pipeline fine-tuning đề xuất}
\label{sec:pipeline-ft}

Vì codebase hiện tại của VNPost chưa có pipeline \ft, phần này trình bày \emph{thiết kế đề
xuất}. Mục tiêu của \ft{} trong bối cảnh này không phải là nạp tri thức (đã do \rag{} đảm
nhiệm) mà là điều chỉnh hành vi: phong cách trả lời, định dạng nghiệp vụ, cách xử lý các
tình huống lặp và mức độ tuân thủ quy tắc.

\subsection{Thu thập và chuẩn bị dữ liệu}

Dữ liệu huấn luyện được tổng hợp từ ba nguồn: (i) log hội thoại thực tế giữa người dùng và
hệ thống, sau khi lọc và ẩn thông tin định danh cá nhân (PII); (ii) các cặp hỏi -- đáp chuẩn
do chuyên gia nghiệp vụ xây dựng; và (iii) dữ liệu tổng hợp do mô hình mạnh sinh ra có kiểm
duyệt. Dữ liệu được định dạng theo mẫu hội thoại (instruction -- response) và chia thành
tập huấn luyện, tập kiểm định và tập kiểm thử.

\subsection{Cấu hình LoRA/QLoRA}

Đề xuất sử dụng LoRA hoặc QLoRA với hạng $r=16$, hệ số $\alpha=32$, dropout nhỏ, và áp dụng
lên các ma trận chiếu trong khối attention và feed-forward (query, key, value, output,
gate, up, down). Mô hình gốc được giữ ở định dạng lượng tử hóa 4 bit khi dùng QLoRA để tiết
kiệm bộ nhớ. Sau huấn luyện, adapter có thể được gộp (merge) vào mô hình gốc hoặc phục vụ
như một adapter riêng trong vLLM.

\begin{figure}[htbp]
\centering
\begin{tikzpicture}[
  node distance=7mm,
  proc/.style={draw=vnpostgreen,fill=boxblue,rounded corners,minimum width=3.8cm,minimum height=9mm,align=center,font=\small},
  arrow/.style={-{Latex[length=2mm]},thick,vnpostdark}
]
\node[proc] (log) {Log hội thoại + cặp hỏi đáp chuẩn + dữ liệu tổng hợp};
\node[proc,below=of log] (clean) {Làm sạch, ẩn PII, chuẩn hóa định dạng};
\node[proc,below=of clean] (split) {Chia train / val / test};
\node[proc,below=of split] (lora) {Huấn luyện LoRA/QLoRA (SFT)};
\node[proc,below=of lora] (eval) {Đánh giá trên val (LLM-as-judge, RAGAS)};
\node[proc,below=of eval] (serve) {Gộp/đóng gói adapter $\rightarrow$ phục vụ vLLM};
\draw[arrow] (log) -- (clean);
\draw[arrow] (clean) -- (split);
\draw[arrow] (split) -- (lora);
\draw[arrow] (lora) -- (eval);
\draw[arrow] (eval) -- (serve);
\end{tikzpicture}
\caption{Pipeline fine-tuning đề xuất cho trợ lý ảo VNPost.}
\label{fig:pipeline-ft}
\end{figure}

\begin{ghichu}
Ghi chú minh bạch: pipeline \ft{} trong báo cáo là \emph{thiết kế đề xuất}; các kết quả liên
quan ở chương \ref{ch:ketqua} (nếu có) được ghi rõ là kết quả minh họa, không phải kết quả
huấn luyện thật trên hạ tầng VNPost.
\end{ghichu}

\section{Kiến trúc lai RAG + Fine-tuning}
\label{sec:kien-truc-lai}

Kiến trúc lai sử dụng một bộ định tuyến để quyết định chiến lược trả lời:

\begin{itemize}
  \item \textbf{Chỉ RAG:} với câu hỏi cần tri thức động hoặc cần trích dẫn nguồn (biểu phí,
  quy định, số hiệu văn bản). Đây là chế độ mặc định cho phần lớn truy vấn.
  \item \textbf{Chỉ mô hình tinh chỉnh:} với câu hỏi về phong cách, định dạng, hướng dẫn quy
  trình ổn định, hoặc khi cần câu trả lời ngắn gọn theo mẫu nghiệp vụ.
  \item \textbf{Kết hợp:} mô hình tinh chỉnh sinh khung câu trả lời theo định dạng nghiệp vụ,
  sau đó \rag{} bổ sung bằng chứng và số liệu cụ thể; mô hình hợp nhất tạo câu trả lời cuối.
\end{itemize}

Bộ định tuyến có thể dựa trên phân loại ý định nhẹ, hoặc để chính tác tử điều phối quyết
định dựa trên prompt. Khi \rag{} không tìm thấy bằng chứng đủ tin cậy, hệ thống chuyển sang
thông báo thiếu thông tin hoặc tạo ticket, tránh sinh câu trả lời không căn cứ.

\begin{figure}[htbp]
\centering
\begin{tikzpicture}[
  node distance=8mm,
  proc/.style={draw=vnpostgreen,fill=boxblue,rounded corners,minimum width=3.2cm,minimum height=9mm,align=center,font=\small},
  dec/.style={draw=vnpostdark,fill=vnpostgreen,text=white,diamond,aspect=2.2,align=center,font=\small,inner sep=1pt},
  arrow/.style={-{Latex[length=2mm]},thick,vnpostdark}
]
\node[proc] (q) {Câu hỏi người dùng};
\node[dec,below=of q] (router) {Bộ định tuyến};
\node[proc,below left=14mm and 26mm of router] (rag) {RAG (tri thức động)};
\node[proc,below right=14mm and 26mm of router] (ft) {Mô hình tinh chỉnh (hành vi)};
\node[proc,below=26mm of router] (comb) {Hợp nhất câu trả lời};
\node[proc,below=of comb] (ans) {Câu trả lời cuối (có trích nguồn)};
\draw[arrow] (q) -- (router);
\draw[arrow] (router) -- (rag);
\draw[arrow] (router) -- (ft);
\draw[arrow] (rag) -- (comb);
\draw[arrow] (ft) -- (comb);
\draw[arrow] (comb) -- (ans);
\end{tikzpicture}
\caption{Kiến trúc lai: bộ định tuyến chọn RAG, mô hình tinh chỉnh hoặc kết hợp.}
\label{fig:kien-truc-lai}
\end{figure}

\section{Hạ tầng và bảo mật}
\label{sec:hatang}

\subsection{Hạ tầng phục vụ}

Toàn bộ hệ thống được triển khai trên Kubernetes để bảo đảm khả năng mở rộng và ổn định.
Các mô hình được phục vụ cục bộ bằng vLLM với lượng tử hóa FP8. Các mô hình chính gồm: mô
hình sinh ngôn ngữ (ví dụ họ Qwen với các biến thể 27B--35B), mô hình nhúng văn bản (ví dụ
Qwen3-Embedding-8B), mô hình xếp hạng lại (Qwen3-Reranker-8B) và mô hình OCR. Qdrant đảm
nhiệm lưu trữ vector.

\subsection{Bảo mật dữ liệu}

Yêu cầu bảo mật được thực thi ở nhiều lớp: (i) mô hình chạy cục bộ, dữ liệu không rời hạ
tầng; (ii) kiểm soát truy cập theo vai trò (RBAC) đối với dữ liệu nhạy cảm; (iii) ẩn PII
trong log trước khi dùng cho phân tích hoặc \ft; và (iv) phân quyền ở tầng dữ liệu, bảo
đảm người dùng chỉ truy cập được phạm vi cho phép.

\section*{Kết chương}
\addcontentsline{toc}{section}{Kết chương}
Chương này đã phân tích bài toán VNPost, xác định yêu cầu và đề xuất kiến trúc đa tác tử
kết hợp \rag{} lai với \ft{} theo mô hình lai. Chương tiếp theo mô tả việc triển khai thực
nghiệm các thành phần này cùng các tham số cụ thể.
